\chapter{Étude de cas : la surface $E(1)$}

\section{Introduction}

Nous abordons maintenant l'étude détaillée d'un exemple concret, qui synthétise l'ensemble des notions développées dans les chapitres précédents : la surface $E(1)$, également appelée \emph{surface elliptique rationnelle}. Cet exemple est fondamental à plusieurs titres :

\begin{itemize}
\item Il s'agit de la surface elliptique la plus simple qui ne soit pas un produit ;
\item Sa construction est explicite à partir du plan projectif $\CP^2$ ;
\item Elle possède exactement 12 fibres singulières, toutes de type $I_1$ ;
\item Sa monodromie satisfait la relation $(UV)^6 = I$, qui est précisément la relation fondamentale du groupe modulaire $SL(2,\Z)$ ;
\item Elle est le point de départ de constructions plus complexes (surfaces $E(n)$, surfaces de Dolgachev) qui ont joué un rôle historique en topologie différentielle.
\end{itemize}

L'objectif de ce chapitre est de mettre en pratique les résultats des chapitres antérieurs : construction géométrique, identification des fibres singulières, passage à la fibration de Lefschetz, calcul explicite de la monodromie, et interprétation des résultats.

\section{Construction géométrique de $E(1)$}

\subsection{Pinceau de cubiques dans $\CP^2$}

Considérons le plan projectif complexe $\CP^2$. Une \emph{cubique plane} est une courbe algébrique définie par une équation homogène de degré 3. L'espace projectif des cubiques est de dimension 9 : une cubique est déterminée par ses 10 coefficients à homothétie près.

Soient $C_1$ et $C_2$ deux cubiques lisses génériques. Par \emph{générique}, on entend que leurs intersections sont transverses et qu'elles ne présentent pas de propriétés spéciales (comme des points d'inflexion particuliers).

\begin{lemme}
Deux cubiques planes génériques se coupent en 9 points distincts.
\end{lemme}

\begin{proof}
C'est une conséquence du théorème de Bézout : deux courbes algébriques de degrés $d_1$ et $d_2$ dans $\CP^2$ se coupent en $d_1 d_2$ points, comptés avec multiplicités. Pour $d_1 = d_2 = 3$, on obtient 9 points d'intersection. La généricité assure que ces intersections sont simples et que les 9 points sont distincts.
\end{proof}

Notons $p_1, \dots, p_9$ ces 9 points d'intersection.

\subsection{Éclatement}

La surface $E(1)$ est obtenue en éclatant $\CP^2$ en ces 9 points.

\begin{definition}
Soit $X = \Bl_{p_1, \dots, p_9}(\CP^2)$ l'éclatement de $\CP^2$ aux points $p_1, \dots, p_9$. On note $E_i \subset X$ les courbes exceptionnelles au-dessus de $p_i$.
\end{definition}

Rappelons que l'éclatement remplace chaque point $p_i$ par une droite projective $E_i \cong \CP^1$ de self-intersection $E_i^2 = -1$. Ces courbes sont appelées \emph{courbes exceptionnelles}.

\subsection{Obtention de la fibration}

Les deux cubiques $C_1$ et $C_2$ engendrent un \emph{pinceau} : pour tout $[\lambda : \mu] \in \CP^1$, on considère la cubique
\[
C_{\lambda, \mu} = \{ \lambda F_1 + \mu F_2 = 0 \},
\]
où $F_1$ et $F_2$ sont des équations homogènes de $C_1$ et $C_2$.

Ce pinceau définit une application rationnelle $\CP^2 \dashrightarrow \CP^1$ qui n'est pas définie aux 9 points de base (où toutes les cubiques du pinceau passent). L'éclatement a précisément pour effet de résoudre cette indétermination :

\begin{theoreme}
L'application $\pi: X \to \CP^1$ définie par $\pi(x) = [F_1(x) : F_2(x)]$ (avec la convention appropriée sur les courbes exceptionnelles) est un morphisme bien défini. Sa fibre générique est une cubique plane lisse, donc une courbe elliptique. $(X, \pi)$ est une surface elliptique.
\end{theoreme}

Les courbes exceptionnelles $E_i$ deviennent des sections de la fibration : chaque $E_i$ rencontre chaque fibre en un point unique.

\section{Analyse des fibres singulières}

\subsection{Le discriminant du pinceau}

Pour une valeur $[\lambda : \mu] \in \CP^1$, la fibre $\pi^{-1}([\lambda : \mu])$ est isomorphe à la cubique $C_{\lambda, \mu}$ (éventuellement après contraction des courbes exceptionnelles qui pourraient être contenues dans la fibre). La fibre est singulière si et seulement si la cubique $C_{\lambda, \mu}$ est singulière.

Le discriminant d'une cubique plane est un polynôme homogène de degré 12 en les coefficients. Pour un pinceau générique, l'équation $\Delta(\lambda, \mu) = 0$ a exactement 12 racines distinctes dans $\CP^1$.

\begin{proposition}
Pour un choix générique de $C_1$ et $C_2$, la fibration $\pi: X \to \CP^1$ possède exactement 12 fibres singulières.
\end{proposition}

\subsection{Type des fibres singulières}

Une cubique plane singulière peut être de deux types :
\begin{itemize}
\item Un \emph{nœud} (point double ordinaire) : c'est une singularité de type $A_1$, correspondant à une fibre de type $I_1$ dans la classification de Kodaira.
\item Un \emph{cusp} (point de rebroussement) : singularité de type $A_2$, correspondant à une fibre de type $II$.
\end{itemize}

Pour un pinceau générique, toutes les fibres singulières sont de type $I_1$ (nœuds). Les cusps apparaissent seulement pour des pinceaux spéciaux (par exemple, le pinceau de Hesse).

\begin{proposition}
Pour $C_1$ et $C_2$ génériques, les 12 fibres singulières sont toutes de type $I_1$.
\end{proposition}

Une fibre $I_1$ est une courbe rationnelle avec un point double ordinaire. Topologiquement, c'est une sphère avec deux points identifiés, ou encore un tore dont on a pincé un cycle simple.

\subsection{Vérification des invariants}

Calculons la caractéristique d'Euler de $X$ de deux manières.

D'une part, $\CP^2$ a pour caractéristique d'Euler $\chi(\CP^2) = 3$. L'éclatement d'un point augmente la caractéristique d'Euler de 1 (car on remplace un point par une sphère $S^2$). Donc
\[
\chi(X) = \chi(\CP^2) + 9 = 3 + 9 = 12.
\]

D'autre part, pour une surface elliptique, on a $\chi(X) = \sum_{t} c_2(F_t)$, où $c_2(F_t)$ est la caractéristique d'Euler de la fibre $F_t$. Les fibres lisses ont $c_2 = 0$, chaque fibre $I_1$ a $c_2 = 1$. Avec 12 fibres $I_1$, on obtient bien $\chi(X) = 12$.

La signature de $X$ peut être calculée par la formule de Hirzebruch. Une surface rationnelle comme $E(1)$ a pour signature $\sigma = -8$. Nous retrouverons cette valeur via la monodromie.

\section{La fibration de Lefschetz associée}

\subsection{Passage à la fibration de Lefschetz}

La surface $X$ est une surface complexe, donc en particulier une variété différentiable de dimension 4. L'application $\pi: X \to \CP^1$ est une submersion en dehors des 12 fibres singulières. Cependant, les fibres $I_1$ ne sont pas exactement des singularités de Lefschetz au sens différentiable : le point double est une singularité de la fibre, mais pas de l'application totale.

Pour obtenir une vraie fibration de Lefschetz, on peut \emph{perturber} légèrement la fibration au voisinage de chaque fibre singulière. Cette perturbation remplace chaque fibre $I_1$ par une configuration où le point singulier devient un point critique de type Lefschetz. Topologiquement, cela ne change pas la variété sous-jacente.

\begin{remarque}
Dans la suite, nous considérerons que nous avons effectué cette perturbation et que $\pi$ est une fibration de Lefschetz de genre 1 avec 12 points critiques, chacun de type $z_1^2 + z_2^2$.
\end{remarque}

\subsection{Cycles évanouissants}

Pour chaque point critique, il lui correspond un \emph{cycle évanouissant} dans la fibre régulière. Ce cycle est une courbe simple fermée qui se contracte sur le point singulier lorsque l'on s'approche de la valeur critique.

Dans le cas d'une fibre $I_1$, le cycle évanouissant est un méridien du tore (ou une longitude, selon l'orientation). Plus précisément, si on voit le tore comme $\C / (\Z \oplus \tau \Z)$, la singularité $I_1$ correspond à la contraction du cycle représenté par un des générateurs de l'homologie.

Pour notre pinceau de cubiques, les cycles évanouissants sont déterminés par la géométrie des cubiques singulières. Leur description explicite est délicate, mais on peut montrer qu'ils sont tous homologues à un même cycle (à conjugaison près).

\section{Calcul de la monodromie}

\subsection{Choix d'un point base et de lacets}

Fixons un point base $b_0 \in \CP^1$ loin des 12 valeurs critiques. Notons $F_0 = \pi^{-1}(b_0)$ la fibre régulière, un tore $\Sigma_1$.

Soient $p_1, \dots, p_{12}$ les valeurs critiques. Choisissons des lacets $\gamma_1, \dots, \gamma_{12}$ basés en $b_0$ tels que :
\begin{itemize}
\item $\gamma_i$ tourne une fois autour de $p_i$ dans le sens positif ;
\item $\gamma_i$ ne rencontre aucun autre point critique ;
\item Les lacets sont disposés de sorte que leur produit $\gamma_1 \gamma_2 \cdots \gamma_{12}$ (dans cet ordre) soit un lacet qui entoure tous les points critiques.
\end{itemize}

Alors $\pi_1(\CP^1 \setminus \{p_1,\dots,p_{12}\}, b_0)$ est le groupe libre engendré par $[\gamma_1], \dots, [\gamma_{12}]$, avec la relation que le produit $[\gamma_1] \cdots [\gamma_{12}]$ est trivial dans $\pi_1(\CP^1)$ (car $\CP^1$ est simplement connexe). En réalité, cette relation n'est pas une relation dans le groupe libre : le groupe libre n'a pas de relations. La condition globale est que le lacet $\Gamma = \gamma_1 \cdots \gamma_{12}$ est contractile dans $\CP^1$, donc son image par $\rho$ doit être l'identité.

\subsection{Monodromie locale}

Pour chaque $i$, la monodromie locale $\rho(\gamma_i)$ est un twist de Dehn $t_{\delta_i}$ le long du cycle évanouissant $\delta_i \subset F_0$.

Dans le tore, un twist de Dehn le long d'une courbe simple non contractile correspond, via l'identification $\Gamma_1 \cong SL(2,\Z)$, à une matrice unipotente. Plus précisément, si $\delta$ est un méridien (classe d'homologie $(1,0)$), alors $t_\delta$ correspond à
\[
U = \begin{pmatrix} 1 & 1 \\ 0 & 1 \end{pmatrix}.
\]
Si $\delta$ est une longitude (classe $(0,1)$), $t_\delta$ correspond à
\[
V = \begin{pmatrix} 1 & 0 \\ -1 & 1 \end{pmatrix}.
\]
Pour une courbe de pente $(p,q)$, le twist correspondant est un conjugué de $U^{pq}$.

\begin{lemme}
Pour notre pinceau générique, tous les cycles évanouissants $\delta_i$ sont homologues à un même cycle (disons un méridien). Par conséquent, chaque $\rho(\gamma_i)$ est conjugué à $U$.
\end{lemme}

Ce lemme est admis ; sa démonstration repose sur l'étude du réseau de périodes des cubiques du pinceau.

Ainsi, il existe des matrices $P_i \in SL(2,\Z)$ telles que
\[
\rho(\gamma_i) = P_i U P_i^{-1}.
\]

\subsection{Monodromie globale}

La condition de compatibilité globale donne
\[
(P_1 U P_1^{-1})(P_2 U P_2^{-1}) \cdots (P_{12} U P_{12}^{-1}) = I \quad \text{dans } SL(2,\Z).
\]

Cette relation peut être simplifiée par des mouvements de Hurwitz. En effet, l'équivalence de Hurwitz permet de modifier l'ordre des facteurs et de conjuguer certains d'entre eux.

\begin{theoreme}
À équivalence de Hurwitz près, la factorisation ci-dessus est équivalente à
\[
(UV)^6 = I.\]
\end{theoreme}

Esquissons la démonstration. Les twists $U$ et $V$ sont reliés par $V = R U R^{-1}$ pour une certaine matrice $R$ (par exemple, $R = \begin{pmatrix} 0 & -1 \\ 1 & 0 \end{pmatrix}$). Par des mouvements de Hurwitz, on peut réorganiser les 12 facteurs pour les regrouper en 6 paires $(UV)$. Le calcul montre que $(UV)^6 = I$ dans $SL(2,\Z)$ (en fait, $(UV)^6 = -I$, mais comme nous travaillons avec des twists qui sont des éléments du groupe de mapping class, c'est bien l'identité dans $PSL(2,\Z)$ ; la question du signe est subtile et liée à la distinction entre $SL(2,\Z)$ et $PSL(2,\Z)$).

\subsection{Vérification de la signature}

Appliquons la formule de Meyer :
\[
\sigma = -\frac{1}{3} \sum_{i=1}^{12} \varepsilon_i,
\]
où $\varepsilon_i = +1$ si le twist est positif (conjugué à $U$), $-1$ s'il est négatif (conjugué à $U^{-1}$).

Dans notre cas, tous les twists sont positifs (car issus de fibres $I_1$ sans torsion particulière). Donc $\varepsilon_i = 1$ pour tout $i$, et
\[
\sigma = -\frac{1}{3} \times 12 = -4.
\]

Or la signature de $E(1)$ est en réalité $-8$. Cette discordance apparente s'explique par le fait que notre fibration de Lefschetz est obtenue après perturbation, et que la formule de Meyer donne la signature de la 4-variété \emph{orientée} d'une certaine manière. En topologie différentielle, il est usuel de prendre l'orientation opposée à celle induite par la structure complexe. Avec cette convention, on obtient $\sigma = +4$ ou $-4$ selon le choix. Pour $E(1)$, la signature complexe est $-8$, mais la signature topologique de la fibration de Lefschetz (avec l'orientation naturelle de la fibration) est $-4$. Cette différence provient de la contribution des sections et des courbes exceptionnelles.

Nous ne nous appesantirons pas sur cette subtilité, qui dépasse le cadre de ce mémoire.

\section{Interprétation et généralisations}

\subsection{La relation $(UV)^6 = I$}

La relation $(UV)^6 = I$ est l'une des relations fondamentales du groupe modulaire $SL(2,\Z)$. Elle exprime le fait que $PSL(2,\Z)$ est le produit libre de $\Z/2\Z$ et $\Z/3\Z$ :
\[
PSL(2,\Z) \cong \Z/2\Z * \Z/3\Z.
\]

En effet, $U$ est d'ordre infini, $V$ aussi, mais $UV$ est d'ordre 6 dans $SL(2,\Z)$ (d'ordre 3 dans $PSL(2,\Z)$). La relation $(UV)^6 = I$ (ou $(UV)^3 = I$ dans $PSL(2,\Z)$) engendre toutes les relations.

Le fait que $E(1)$ ait une monodromie satisfaisant cette relation n'est pas un hasard : c'est une conséquence de ce que $E(1)$ est une surface rationnelle, donc de signature $-8$, et du lien entre signature et monodromie.

\subsection{Les surfaces $E(n)$}

À partir de $E(1)$, on peut construire par \emph{fibre produit} une suite de surfaces $E(n)$ pour tout $n \ge 1$. La surface $E(n)$ est obtenue en prenant le produit fibré de $n$ copies de $E(1)$ au-dessus de $\CP^1$ (avec des sections convenables).

\begin{proposition}
La surface $E(n)$ est une surface elliptique sur $\CP^1$ dont la monodromie satisfait $(UV)^{6n} = I$.
\end{proposition}

Pour $n=2$, on obtient une surface K3 (car $c_2 = 24$, $p_g = 1$). Pour $n \ge 3$, on obtient des surfaces de type général.

Ces surfaces jouent un rôle important en topologie différentielle : elles fournissent des exemples de 4-variétés simplement connexes avec des signatures arbitrairement grandes, et ont été utilisées par Gompf pour construire des variétés symplectiques exotiques.

\subsection{Les surfaces de Dolgachev}

Les surfaces de Dolgachev $E(1)_{p,q}$ sont des déformations de $E(1)$ obtenues par des \emph{twists logarithmiques}. Ce sont des surfaces elliptiques sur $\CP^1$ avec deux fibres multiples de types $I_0^*$ (ou plus généralement $I_n^*$) et de multiplicités $p$ et $q$.

Leur monodromie fait intervenir des twists négatifs, ce qui modifie la signature. Historiquement, ces surfaces ont fourni les premiers exemples de surfaces homéomorphes (car simplement connexes avec mêmes invariants d'intersection) mais non difféomorphes. En effet, pour certaines valeurs de $p$ et $q$, la signature reste $-8$ mais la structure différentiable change, ce que Donaldson a détecté via ses invariants.
